% ============================================================
% preambulo-contrato-comun.tex
% ------------------------------------------------------------
% CONTRATO COMÚN VERSIÓN 1
% ------------------------------------------------------------
% v1 (07/11/2026): SALE DE preambulo-contrato.tex, QUE ERA DE
% LIBROS. ES LO QUE EMITEN LAS DOS HOJAS DE LaTeX,
% docbook-to-latex.xsl (LIBROS) Y jats-to-latex.xsl (REVISTAS), Y
% NO DEPENDE DE LA CLASE NI DEL CAPÍTULO: CORTE Y NOTA EN
% TÍTULOS, TACHADO Y SUBRAYADO, LOS PAQUETES DEL CONTRATO, babel,
% csquotes, CIFRAS, CUADROS, BLOQUES (CITA, EPÍGRAFE, RECUADROS,
% VERSO, CONVERSACIÓN), LLAMADAS DE ATENCIÓN Y CÓDIGO. EL TEXTO
% DE CADA APARTADO SE MUDÓ SIN CAMBIOS, SALVO EL CONTADOR
% gbcodigo (VER «8. CÓDIGO»).
% LO CARGA, PRIMERO, EL CONTRATO DE CADA TIPO DE PROYECTO:
%   libro    preambulo-contrato.tex
%   revista  preambulo-revista-contrato.tex
% UN CAMBIO ACÁ LLEGA A LOS DOS: SE PRUEBA CON UN LIBRO Y CON
% UNA REVISTA ANTES DE SUBIR LA VERSIÓN.
% ------------------------------------------------------------
% DEPENDE DE LAS DEFINICIONES DEL PROYECTO, QUE VAN ANTES:
%   \gbidiomas          LOS IDIOMAS DE babel
%   \ifgbcodigocolor    EL CÓDIGO EN COLOR O EN GRISES (SI NO ESTÁ
%                       DEFINIDO, EN COLOR)
% ============================================================

\def\gbContratoComunVersion{1}

% ============================================================
% CORTE DE LÍNEA EN TÍTULOS (CONTRATO 6, SC-28)
% ------------------------------------------------------------
% El editor marca en la base dónde cortar un título o un
% subtítulo. El generador lo traduce a \gbCorteTitulo SOLO en
% la versión compuesta del título: la que va a la apertura de
% la pieza o de la parte y a la portada. El sumario, el folio y
% los marcadores reciben siempre la versión plana, sin la
% macro: por eso las macros de título de este contrato llevan
% las dos.
%
% \gbtitulo Y \gbsubtitulo VAN COMPUESTOS Y TAMBIÉN LLEGAN A
% pdftitle Y pdfsubject. AHÍ hyperref PASA EL VALOR A CADENA
% PDF, Y UN \\ NO TIENE LUGAR: LA DESACTIVACIÓN LO CAMBIA POR
% UN ESPACIO. VALE ESPACIO Y NO NADA PORQUE EL GENERADOR ESCRIBE
% «Uno\gbCorteTitulo dos», Y EL ESPACIO QUE SIGUE A LA MACRO LO
% CONSUME TeX AL LEER EL NOMBRE.
%
% hyperref SE CARGA EN LA ÚLTIMA CAPA, Y \pdfstringdefDisableCommands
% NO EXISTE ANTES: SE REGISTRA EN EL GANCHO DE CARGA DEL PAQUETE,
% QUE CORRE APENAS hyperref TERMINA DE CARGARSE Y ANTES DEL
% \hypersetup{pdftitle=...} DE ESA MISMA CAPA. EN \AtBeginDocument
% SERÍA TARDE: pdftitle YA ESTARÍA CONVERTIDO.
% ============================================================
\newcommand{\gbCorteTitulo}{\\}
\AddToHook{package/hyperref/after}{%
  \pdfstringdefDisableCommands{\def\gbCorteTitulo{ }}}

% ============================================================
% NOTA AL PIE EN UN TÍTULO DE SECCIÓN (CONTRATO 7, SC-29)
% ------------------------------------------------------------
% \gbNotaTitulo{TEXTO}: LA EMITE docbook-to-latex.xsl EN EL
% ARGUMENTO OBLIGATORIO DE \section Y SUS NIVELES, NUNCA EN EL
% OPCIONAL: EL SUMARIO, EL FOLIO Y LOS MARCADORES VAN SIN NOTA.
%
% CON LA OPCIÓN calcwidth, titlesec COMPONE EL TÍTULO DOS VECES,
% UNA PARA MEDIRLO, Y UN \footnote DIRECTO AVANZA DOS VECES EL
% CONTADOR: LAS NOTAS SALEN 2, 3, 5 (VERIFICADO, GV-70). EL
% MANUAL DE titlesec (§2.9) DA \iftitlemeasuring{MIDE}{COMPONE}
% PARA ESTE CASO: AL MEDIR NO SE EMITE NADA.
%
% SI LA ESTÉTICA NO CARGA titlesec, \iftitlemeasuring NO EXISTE:
% SE PROVEE AL EMPEZAR EL DOCUMENTO, CUANDO LA ESTÉTICA YA SE
% CARGÓ, Y SIEMPRE COMPONE. \providecommand NO TOCA LA DE
% titlesec.
%
% LOS PÁRRAFOS DE LA NOTA LLEGAN SEPARADOS CON \endgraf: EL
% ARGUMENTO DE UN TÍTULO NO ADMITE \par.
% ============================================================
\newcommand{\gbNotaTitulo}[1]{%
  \iftitlemeasuring{}{\footnote{#1}}}
\AtBeginDocument{%
  \providecommand{\iftitlemeasuring}[2]{#2}}

% ============================================================
% TACHADO Y SUBRAYADO (CONTRATO 8, SC-30)
% ------------------------------------------------------------
% \gbTachado{TEXTO} Y \gbSubrayado{TEXTO}: LOS EMITE
% docbook-to-latex.xsl PARA ~~x~~ Y [x]{.underline}.
%
% ulem Y NO \underline: \underline ARMA UNA CAJA Y NO CORTA LA
% LÍNEA, Y LaTeX NO TIENE TACHADO PROPIO. [normalem] ES
% OBLIGATORIO: SIN ÉL, ulem CAMBIA \emph POR SUBRAYADO EN TODO
% EL LIBRO.
%
% \DeclareRobustCommand PORQUE PUEDEN QUEDAR EN UN TÍTULO DE
% SECCIÓN, QUE VIAJA AL SUMARIO, Y LOS COMANDOS DE ulem SON
% FRÁGILES. EN LOS MARCADORES DEL PDF SE QUEDA EL TEXTO SOLO.
% ============================================================
\usepackage[normalem]{ulem}
\DeclareRobustCommand{\gbTachado}[1]{\sout{#1}}
\DeclareRobustCommand{\gbSubrayado}[1]{\uline{#1}}
\AddToHook{package/hyperref/after}{%
  \pdfstringdefDisableCommands{%
    \def\gbTachado#1{#1}%
    \def\gbSubrayado#1{#1}}}

% ============================================================
% 1. PAQUETES QUE EL CONTRATO NECESITA
% ------------------------------------------------------------
% SOLO LO QUE HACE FALTA PARA QUE UN FRAGMENTO COMPILE. LA
% ESTÉTICA CARGA LO SUYO APARTE. TODOS ESTÁN EN TeX Live Y
% NINGUNO PIDE --shell-escape.
% ============================================================
\usepackage{etoolbox}
\usepackage{graphicx}
\usepackage{booktabs}
\usepackage{array}
% fancyvrb: EL Verbatim CON commandchars DEL CÓDIGO COLOREADO (SC-42)
\usepackage{fancyvrb}
\usepackage{ragged2e}
\usepackage{xcolor}
\usepackage[most]{tcolorbox}
% SEPARADOR ORNAMENTAL \froufrou (SC-36). CARGA tikz Y fourier-orns
\usepackage{froufrou}
% VERSO (SC-49): REDEFINE verse Y TRAE \vin, \\! Y \stanzaskip. EN
% TeX Live 2023 (MINT) ES LA 2.4a; LO QUE SE USA EXISTE DESDE LA 2.2
\usepackage{verse}

% ------------------------------------------------------------
% IDIOMA
% ------------------------------------------------------------
% babel VA EN EL CONTRATO Y NO EN LA ESTÉTICA, POR DOS RAZONES.
% LA PRIMERA ES DE ORDEN: biblatex PIDE CARGARSE DESPUÉS DE
% babel Y csquotes, Y LA ESTÉTICA VA DESPUÉS DE ESTA CAPA.
% LA SEGUNDA ES DE FONDO: EL IDIOMA DECIDE LAS COMILLAS QUE
% PRODUCE \enquote, LA PARTICIÓN DE PALABRAS Y LAS CADENAS DE
% LA BIBLIOGRAFÍA. SON GARANTÍAS DEL CONTRATO, NO DECISIONES
% DE DISEÑO: UNA EDITORIAL QUE REEMPLAZA LA ESTÉTICA NO DEBE
% PODER CAMBIAR EL IDIOMA DEL LIBRO.
\usepackage[\gbidiomas]{babel}
\frenchspacing

% COMILLAS SEGÚN EL IDIOMA ACTIVO: EN ESPAÑOL DA « ».
% EL XSLT EMITE \enquote Y NUNCA COMILLAS LITERALES.
\usepackage[autostyle=true]{csquotes}

% ------------------------------------------------------------
% CIFRAS Y CANTIDADES (CONTRATO 15, SC-64, SC-65)
% ------------------------------------------------------------
% LA NORMA INTERNACIONAL (SI Brochure §5.4): PUNTO DECIMAL,
% GRUPOS DE A TRES CON UN ESPACIO DESDE CINCO DÍGITOS, LA UNIDAD
% Y EL PORCENTAJE SEPARADOS DE LA CIFRA, EL GRADO PEGADO.
% detect-all NO SE QUITA: SIN ELLA siunitx COMPONE LA CIFRA EN
% MODO MATEMÁTICO, DONDE babel-spanish PASA EL PUNTO A COMA, Y
% LOS DECIMALES SALEN «12,5» SIN AVISO (GV-123). EN siunitx 3 ES
% UNA OPCIÓN DEPRECADA: DEJA UN AVISO INFORMATIVO EN EL LOG.
% per-mode=symbol: «48 km/h» EN UNA LÍNEA; LA FRACCIÓN VERTICAL
% AGRANDABA LA LÍNEA (DECISIÓN DE ALBERTO).
% EL FORMATEADOR DE cifras.lua APLICA LAS MISMAS REGLAS PARA EL
% HTML Y EL EPUB: SI CAMBIA ESTE \sisetup, CAMBIA EL FILTRO.
\usepackage{siunitx}
\sisetup{output-decimal-marker={.},
  group-separator={\hspace{0.15em}},
  group-minimum-digits=5,
  table-text-alignment=center,
  detect-all,
  per-mode=symbol}
% siunitx NO DEFINE EL EURO (GV-123). LA FUENTE TIENE QUE TENER €
\DeclareSIUnit\euro{€}

% ------------------------------------------------------------
% CUADROS (CONTRATO 16, SC-66)
% ------------------------------------------------------------
% EL ESTILO DE LA CASA: LETRA SIN SERIFA, FILETE ENTRE TODAS LAS
% FILAS Y ENTRE TODAS LAS COLUMNAS, ENCABEZADO EN NEGRITA, PIE
% DEBAJO. LOS FILETES Y LA NEGRITA LOS ESCRIBE EL XSLT; ACÁ, LO
% QUE UNA EDITORIAL PUEDE CAMBIAR SIN TOCARLO. EL CUERPO Y LA
% SEPARACIÓN DE COLUMNAS LOS ELIGE EL MAQUETADOR POR CUADRO
% (cuerpo: normal|chico, separacion: normal|estrecha). LA LETRA
% Y EL CUERPO VALEN PARA TODO EL FLOTANTE, EL PIE INCLUIDO.
% GEMELAS: m_XML.ObtenerPreambuloEmbebido (REVISTAS).
\newcommand{\gbCuadroLetra}{\sffamily}
\newcommand{\gbCuadroCuerpoNormal}{\footnotesize}
\newcommand{\gbCuadroCuerpoChico}{\scriptsize}
\newcommand{\gbCuadroSepNormal}{4pt}
\newcommand{\gbCuadroSepEstrecha}{2pt}
% UN CUADRO CORTADO (cortes:) VA EN UN FLOTANTE POR TRAMO: EL PRIMERO
% TERMINA CON \gbCuadroContinua, A LA DERECHA; LOS QUE SIGUEN EMPIEZAN
% CON \gbCuadroViene, A LA IZQUIERDA, Y REPITEN EL ENCABEZADO. EL PIE,
% SOLO EN EL ÚLTIMO. EL XSLT LOS PONE EN BASTARDILLA.
\newcommand{\gbCuadroContinua}{Continúa en la página siguiente}
\newcommand{\gbCuadroViene}{Viene de la página anterior}
% EL CUADRO APAISADO (rotar: sí): sidewaystable, UNA PÁGINA POR TRAMO
\usepackage{rotating}

% ============================================================
% 6. BLOQUES
% ============================================================

% CITA EN BLOQUE CON ATRIBUCIÓN OPCIONAL.
% NO ES UN EPÍGRAFE: DOCBOOK LOS DISTINGUE Y EL XSLT TAMBIÉN.
\newenvironment{gbcita}
  {\begin{quote}\small}
  {\end{quote}}
\newcommand{\gbatribucion}[1]{\par\vspace{2pt}\hfill\textup{\small#1}\par}

% EPÍGRAFE (SC-35). BLOQUE A LA DERECHA, 60 % DE LA COLUMNA, \small,
% SIN CORTE DE PALABRA. #1 A LA IZQUIERDA; FILETE DE 0,6 PT; #2 A LA
% DERECHA. CON #2 VACÍO NO HAY FILETE. ENTRE PÁRRAFOS DE #1, MEDIA LÍNEA. #1 PUEDE TENER VARIOS PÁRRAFOS
% (\newcommand SIN ASTERISCO ADMITE \par). #2 PUEDE SER UNA CITA: EL
% XSLT RESUELVE UN <biblioref> A \textcite.
% GEMELA DE LA DEFINICIÓN DE m_XML.ObtenerPreambuloEmbebido (REVISTAS):
% UN CAMBIO EN UNA VA EN LAS DOS.
\newcommand{\gbepigrafe}[2]{%
  \par\vspace{.5\baselineskip}%
  {\noindent\hfill
   \begin{minipage}{.6\linewidth}%
     \small\parindent=0pt\parskip=.5em
     \hyphenpenalty=10000\exhyphenpenalty=10000
     \raggedright #1\par
     \if\relax\detokenize{#2}\relax\else
       \parskip=0pt\vspace{1ex}\hrule height 0.6pt\vspace{1ex}%
       {\raggedleft #2\par}%
     \fi
   \end{minipage}\par}%
  \vspace{.5\baselineskip}}

% RECUADRO LATERAL. UN ARGUMENTO: info | supplementary.
\newenvironment{gbsidebar}[1]
  {\begin{tcolorbox}[colback=black!4,colframe=black!4,
     boxrule=0pt,arc=2pt,left=6pt,right=6pt,top=4pt,bottom=4pt,
     before skip=6pt,after skip=6pt]\small\sffamily}
  {\end{tcolorbox}}

% ============================================================
% RECUADROS (CONTRATO 10, SC-41)
% ------------------------------------------------------------
% gbRecuadro: FONDO GRIS AL 7 %, FILETE DE 1 PT Y ESQUINAS DE
% 5 PT, LOS DOS EN gbcolorbase.
% gbRecuadroBarra{BARRA}: FONDO GRIS AL 7 %, SIN FILETE Y CON
% ESQUINAS RECTAS; ARRIBA, UNA BARRA DE UNA PICA EN gbcolorbase
% CON EL TEXTO EN SANS NEGRITA Y MAYÚSCULAS, BLANCO Y CENTRADO
% EN ALTURA: EL MISMO RÓTULO QUE EL TIPO DE ARTÍCULO DE LA
% CABECERA. \MakeUppercase Y NO upper-case() EN EL XSLT: CONSERVA
% LA BASTARDILLA Y DEJA INTACTA LA CLAVE DE UN \autocite (MEDIDO).
% EL TEXTO ES EL DEL CUERPO: breakable PARA QUE UN RECUADRO
% LARGO CRUCE DE PÁGINA; parbox=false PARA QUE LOS PÁRRAFOS
% TENGAN LA SANGRÍA DEL CUERPO (EN MODO MINIPÁGINA QUEDARÍAN
% PEGADOS Y SIN SANGRÍA); EL PRIMERO VA SIN SANGRÍA.
% SIN NOTAS AL PIE: EN UN tcolorbox QUEDAN ADENTRO, CON LETRA
% (MEDIDO); EL FILTRO recuadros.lua LAS FRENA.
% EL COLOR LO DEFINE LA ESTÉTICA (\definecolor{gbcolorbase});
% ESTE \providecolor ES EL RESPALDO DE UNA ESTÉTICA QUE NO LO
% DECLARE. GEMELOS EN m_XML.ObtenerPreambuloEmbebido (REVISTAS):
% UN CAMBIO EN UNO VA EN LOS DOS.
% ============================================================
\providecolor{gbcolorbase}{HTML}{2D5A8E}
\tcbset{
  gbrecuadro/.style={breakable, parbox=false,
    colback=black!7, colframe=gbcolorbase, boxrule=1pt, arc=5pt,
    boxsep=0pt, left=8pt, right=8pt, top=8pt, bottom=8pt,
    before skip=\baselineskip, after skip=\baselineskip,
    before upper=\noindent},
  gbrecuadrobarra/.style={breakable, parbox=false, sharp corners,
    colback=black!7, colframe=black!7, boxrule=0pt,
    colbacktitle=gbcolorbase, coltitle=white, titlerule=0pt,
    boxsep=0pt, left=8pt, right=8pt, top=8pt, bottom=8pt,
    toptitle=0pt, bottomtitle=0pt, lefttitle=8pt, righttitle=8pt,
    before skip=\baselineskip, after skip=\baselineskip,
    before upper=\noindent}}
\newenvironment{gbRecuadro}
  {\begin{tcolorbox}[gbrecuadro]}
  {\end{tcolorbox}}
\newenvironment{gbRecuadroBarra}[1]
  {\begin{tcolorbox}[gbrecuadrobarra,
     title={\parbox[c][1pc][c]{\linewidth}{\sffamily\bfseries\MakeUppercase{#1}}}]}
  {\end{tcolorbox}}

% ============================================================
% VERSO (CONTRATO 14, SC-49)
% ------------------------------------------------------------
% EL ENTORNO verse DEL PAQUETE verse. EL XSLT TERMINA CADA VERSO
% EN \\ Y EL ÚLTIMO DE CADA ESTROFA EN \\!, QUE DEJA \stanzaskip
% (0,75 DE LÍNEA) ENTRE ESTROFAS; EL ÚLTIMO DEL POEMA VA SIN NADA.
% LA SANGRÍA ES \vin (\vgap, 1,5 em) UNA VEZ POR NIVEL; UN VERSO
% QUE NO ENTRA SIGUE A \vindent (3 em). CADA VERSO EMPIEZA CON
% \relax, PARA QUE UN [ O UN * DEL TEXTO NO SE LEA COMO ARGUMENTO
% DE \\ (GV-93).
% gbverso: EL POEMA, EN BASTARDILLA, A \leftmargini DEL MARGEN.
% gbversoepigrafe: EL VERSO DE UN EPÍGRAFE, DENTRO DEL PRIMER
% ARGUMENTO DE \gbepigrafe: LA LETRA DEL EPÍGRAFE, SIN MARGEN
% PROPIO Y SIN EL \topsep DE LA LISTA, QUE SEPARARÍA EL ÚLTIMO
% VERSO DEL FILETE. LOS DOS SALEN DE \@listi, QUE \list LEE AL
% ABRIR: SE REDEFINE LOCAL AL ENTORNO. MEDIDO: DEL CUERPO AL
% PRIMER VERSO Y DEL ÚLTIMO VERSO A LA ATRIBUCIÓN, LA MISMA
% DISTANCIA QUE EN UN EPÍGRAFE EN PROSA (GV-94). \parsep E
% \itemsep LOS PONE verse DESPUÉS (\stanzaskip ENTRE ESTROFAS).
% GEMELOS EN m_XML.ObtenerPreambuloEmbebido (REVISTAS): UN CAMBIO
% EN UNOS VA EN LOS OTROS.
% ============================================================
\newenvironment{gbverso}
  {\begin{verse}\itshape}
  {\end{verse}}
\makeatletter
\newenvironment{gbversoepigrafe}
  {\def\@listi{\leftmargin\z@\topsep\z@\partopsep\z@
     \itemsep\z@\parsep\z@}%
   \begin{verse}}
  {\end{verse}}
\makeatother

% ============================================================
% CONVERSACIÓN (CONTRATO 12, SC-43)
% ------------------------------------------------------------
% gbconversacion: EL qandaset role="conversacion" ENTERO. SANS,
% \small; TODO EL BLOQUE CORRIDO 14 PT A LA IZQUIERDA (\leftskip), EL
% MISMO MARGEN QUE LA CITA (quote, EN preambulo-estetica.tex), SIN
% MARGEN A LA DERECHA; SIN SANGRÍA DE PRIMERA LÍNEA. \leftskip Y NO
% UNA LISTA COMO LA CITA: DENTRO DE UNA LISTA EL ESPACIO ENTRE PÁRRAFOS
% ES \parsep, NO \parskip. LAS NOTAS NO LO HEREDAN: \@parboxrestore
% LO PONE EN CERO. UN ESPACIO CHICO (\smallskipamount) ENTRE
% TURNOS Y ENTRE LOS PÁRRAFOS DE UN TURNO, QUE SIN SANGRÍA ES LO
% ÚNICO QUE LOS SEPARA; \medskip AL ABRIR Y AL CERRAR, QUE DA LA
% SENSACIÓN DE BLOQUE.
% \gbetiqueta{ETIQUETA}: ABRE EL PRIMER PÁRRAFO DEL TURNO, EN
% LÍNEA. LA HOJA LA ESCRIBE YA EN MAYÚSCULAS; PESO NORMAL.
% LA PREGUNTA NO QUEDA SOLA AL PIE: LA HOJA ESCRIBE \nopagebreak[4]
% DESPUÉS DE CADA UNA.
% GEMELO DE LA DEFINICIÓN DE m_XML.ObtenerPreambuloEmbebido
% (REVISTAS): UN CAMBIO EN UNA VA EN LAS DOS.
% ============================================================
\newenvironment{gbconversacion}
  {\par\medskip
   \begingroup\sffamily\small
   \setlength{\leftskip}{14pt}%
   \setlength{\parindent}{0pt}%
   \setlength{\parskip}{\smallskipamount}}
  {\par\endgroup\medskip}
\newcommand{\gbetiqueta}[1]{#1\enspace}

% ============================================================
% 7. LLAMADAS DE ATENCIÓN
% ============================================================
% CADA LLAMADA DE ATENCIÓN ES UN tcolorbox CON TÍTULO FIJO.
% NO SE USA UN .style DE pgfkeys PARAMETRIZADO: EL #1 DE UN
% \newcommand DENTRO DE UN .style LO CONSUME pgfkeys Y DA
% "You already have nine parameters", QUE NO DICE NADA SOBRE
% LA CAUSA. CON UN ENTORNO POR CLASE EL PROBLEMA NO EXISTE.
\tcbset{gbadmon/.style={
  colback=black!3, colframe=black!25,
  boxrule=0.4pt, arc=1pt,
  left=6pt, right=6pt, top=4pt, bottom=4pt,
  before skip=6pt, after skip=6pt,
  fonttitle=\bfseries\sffamily\small}}

\newenvironment{gbnota}
  {\begin{tcolorbox}[gbadmon,title={Nota}]\small}{\end{tcolorbox}}
\newenvironment{gbconsejo}
  {\begin{tcolorbox}[gbadmon,title={Sugerencia}]\small}{\end{tcolorbox}}
\newenvironment{gbadvertencia}
  {\begin{tcolorbox}[gbadmon,title={Advertencia}]\small}{\end{tcolorbox}}
\newenvironment{gbimportante}
  {\begin{tcolorbox}[gbadmon,title={Importante}]\small}{\end{tcolorbox}}
\newenvironment{gbprecaucion}
  {\begin{tcolorbox}[gbadmon,title={Precaución}]\small}{\end{tcolorbox}}

% ============================================================
% 8. CÓDIGO (CONTRATO 11, SC-42)
% ------------------------------------------------------------
% EL COLOR NO LO PONE LaTeX: LO PONE EL RESALTADOR DE PANDOC
% (skylighting) EN UN PASO PREVIO (colorear_codigo.sh), QUE
% ESCRIBE CADA LÍNEA CON LAS MACROS DE TOKEN (\KeywordTok{for}
% Y LAS DEMÁS). ACÁ SOLO SE DEFINE CÓMO SE VE CADA UNA. NI
% listings NI minted: listings NO CONOCE JavaScript, JSON, YAML
% NI Lua (MEDIDO), Y minted EXIGE --shell-escape Y Pygments.
%
% gbCodigo{RÓTULO}: CAJA OSCURA CON CABECERA, EL RÓTULO DEL
% LENGUAJE ARRIBA Y LAS LÍNEAS EN UN Verbatim CON commandchars.
% breakable: UN BLOQUE LARGO CRUZA DE PÁGINA.
% \gbnl{N}: NÚMERO DE LÍNEA, AL MARGEN Y EN GRIS.
% \gbnc: CONTINUACIÓN DE UN CORTE MANUAL, SIN NÚMERO.
% \gbret: EL ↩ DEL CORTE MANUAL, EN EL GRIS DE LOS NÚMEROS.
% \gbPieCodigo{ID}{TEXTO}: EL PIE DE UN LISTADO, DEBAJO DEL
% BLOQUE COMO EN LAS FIGURAS: «Código 2.3.», CON EL CONTADOR
% gbcodigo (DECLARADO CON LOS DE LA NUMERACIÓN POR PIEZA).
%
% \ifgbcodigocolor: VERDADERO, PALETA Gruvbox OSCURA (EL TEMA
% DE LA APLICACIÓN); FALSO, FONDO BLANCO, TEXTO NEGRO Y SIN
% COLOR, PARA IMPRENTA. LA BANDERA LA ESCRIBE EL GENERADOR EN
% LAS PRESTACIONES DE SALIDA (estado «publicado»); SI NO LLEGA,
% VA EL COLOR. LA NEGRITA Y LA BASTARDILLA DE LOS TOKENS SE
% CONSERVAN EN LOS DOS CASOS.
%
% LA PALETA ES LA DE .hidden/themes/gruvbox.theme. GEMELOS: EL
% CSS DEL CÓDIGO EN gbpublisher.css Y EN jats-to-html.xsl, Y
% ESTA MISMA SECCIÓN EN m_XML.ObtenerPreambuloEmbebido
% (REVISTAS). UN CAMBIO DE COLOR VA EN LOS CUATRO.
% ============================================================
\expandafter\ifx\csname ifgbcodigocolor\endcsname\relax
  \expandafter\newif\csname ifgbcodigocolor\endcsname
  \gbcodigocolortrue
\fi

\ifgbcodigocolor
  \definecolor{gbcodfondo}{HTML}{282828}
  \definecolor{gbcodmarco}{HTML}{282828}
  \definecolor{gbcodtexto}{HTML}{EBDBB2}
  \definecolor{gbcodcabeza}{HTML}{3C3836}
  \definecolor{gbcodrotulo}{HTML}{A89984}
  \definecolor{gbcodnumero}{HTML}{7C6F64}
\else
  \definecolor{gbcodfondo}{gray}{1}
  \definecolor{gbcodmarco}{gray}{0}
  \definecolor{gbcodtexto}{gray}{0}
  \definecolor{gbcodcabeza}{gray}{0.9}
  \definecolor{gbcodrotulo}{gray}{0}
  \definecolor{gbcodnumero}{gray}{0.45}
\fi

% \gbDefTok{NOMBRE}{COLOR HTML}{ESTILO}: DEFINE \NOMBRETok. EL
% ESTILO ES \textbf, \textit O \relax. SIN COLOR, SOLO EL ESTILO.
\newcommand{\gbDefTok}[3]{%
  \ifgbcodigocolor
    \expandafter\newcommand\csname #1Tok\endcsname[1]{\textcolor[HTML]{#2}{#3{##1}}}%
  \else
    \expandafter\newcommand\csname #1Tok\endcsname[1]{#3{##1}}%
  \fi}
% LAS 31 CLASES DE skylighting
\gbDefTok{Normal}{EBDBB2}{\relax}
\gbDefTok{Variable}{EBDBB2}{\relax}
\gbDefTok{Keyword}{83A598}{\textbf}
\gbDefTok{ControlFlow}{83A598}{\textbf}
\gbDefTok{DataType}{FABD2F}{\relax}
\gbDefTok{Attribute}{FABD2F}{\relax}
\gbDefTok{DecVal}{D3869B}{\relax}
\gbDefTok{BaseN}{D3869B}{\relax}
\gbDefTok{Float}{D3869B}{\relax}
\gbDefTok{Constant}{D3869B}{\relax}
\gbDefTok{Char}{B8BB26}{\relax}
\gbDefTok{String}{B8BB26}{\relax}
\gbDefTok{VerbatimString}{B8BB26}{\relax}
\gbDefTok{SpecialString}{B8BB26}{\relax}
\gbDefTok{SpecialChar}{FE8019}{\relax}
\gbDefTok{Import}{FE8019}{\relax}
\gbDefTok{Preprocessor}{FE8019}{\relax}
\gbDefTok{Comment}{928374}{\textit}
\gbDefTok{Documentation}{928374}{\textit}
\gbDefTok{Annotation}{928374}{\textit}
\gbDefTok{CommentVar}{928374}{\textit}
\gbDefTok{Information}{928374}{\textit}
\gbDefTok{RegionMarker}{928374}{\relax}
\gbDefTok{Other}{8EC07C}{\relax}
\gbDefTok{Function}{8EC07C}{\relax}
\gbDefTok{Operator}{8EC07C}{\relax}
\gbDefTok{BuiltIn}{8EC07C}{\relax}
\gbDefTok{Extension}{8EC07C}{\relax}
\gbDefTok{Warning}{FE8019}{\textbf}
\gbDefTok{Alert}{FB4934}{\textbf}
\gbDefTok{Error}{FB4934}{\relax}

\newcommand{\gbnl}[1]{\llap{\textcolor{gbcodnumero}{#1}\hspace{1.2em}}}
\newcommand{\gbnc}{\relax}
\newcommand{\gbret}{\textcolor{gbcodnumero}{\symbol{"21A9}}}

\tcbset{gbcodigo/.style={breakable, colback=gbcodfondo,
  colframe=gbcodmarco, coltext=gbcodtexto, boxrule=0.4pt, arc=3pt,
  colbacktitle=gbcodcabeza, coltitle=gbcodrotulo, titlerule=0pt,
  fonttitle=\ttfamily\scriptsize, toptitle=1pt, bottomtitle=1pt,
  boxsep=0pt, left=3em, right=6pt, top=5pt, bottom=5pt,
  lefttitle=6pt, before skip=\baselineskip, after skip=\baselineskip}}
% \VerbatimEnvironment: fancyvrb BUSCA \end{gbCodigo} Y NO
% \end{Verbatim}. ES LA FORMA DOCUMENTADA DE ENVOLVER UN Verbatim.
\newenvironment{gbCodigo}[1]
  {\VerbatimEnvironment
   \begin{tcolorbox}[gbcodigo,title={#1}]%
   \begin{Verbatim}[commandchars=\\\{\},fontsize=\footnotesize]}
  {\end{Verbatim}\end{tcolorbox}}

% EL CONTADOR DE LOS LISTADOS (SC-42). ANTES SE DECLARABA EN EL
% CONTRATO DE LIBROS CON \newcounter{gbcodigo}[chapter]; LA CLASE
% article NO TIENE chapter. EL REINICIO POR CAPÍTULO LO AGREGA EL
% CONTRATO DE LIBROS CON \counterwithin*, QUE HACE LO MISMO QUE EL
% ARGUMENTO OPCIONAL (\@addtoreset) Y NO TOCA \thegbcodigo
\newcounter{gbcodigo}

% EL PIE VA PEGADO AL BLOQUE: SE DESCUENTA PARTE DEL after skip
\newcommand{\gbPieCodigo}[2]{%
  \par\nopagebreak\vspace{-.5\baselineskip}%
  \refstepcounter{gbcodigo}\label{#1}%
  {\small\noindent\textbf{Código~\thegbcodigo.}\space #2\par}%
  \vspace{.5\baselineskip}}
